% Figure: green-hydrogen power-to-power efficiency cascade. Each bar is the
% energy remaining after a stage, normalised to 100 units of input
% electricity, so the bar height is the cumulative efficiency. The last bar
% is the electricity recovered through the fuel cell; the gaps are the
% stage losses. Values illustrative, current as of 2024.
\begin{figure}[htbp]
\centering
\begin{tikzpicture}
\begin{axis}[
  width=13cm, height=7.5cm,
  ylabel={energy remaining (\% of input electricity)},
  symbolic x coords={input, electrolysis, compression, storage, fuelcell},
  xtick={input, electrolysis, compression, storage, fuelcell},
  xticklabels={input\\electricity, after\\electrolysis, after\\compression, after\\storage, after\\fuel cell},
  xticklabel style={align=center, font=\scriptsize},
  ymin=0, ymax=105,
  ytick={0,20,40,60,80,100},
  ybar,
  bar width=20pt,
  nodes near coords,
  nodes near coords style={font=\scriptsize},
  tick label style={font=\scriptsize},
  label style={font=\small},
  enlarge x limits=0.12,
]
\addplot[draw=engblue, fill=engblue!55] coordinates {
  (input,100)
  (electrolysis,65)
  (compression,59)
  (storage,56)
  (fuelcell,31)
};
\end{axis}
\end{tikzpicture}
\caption[Green-hydrogen power-to-power efficiency cascade]{The
power-to-power efficiency cascade for green hydrogen, normalised to $100$
units of input electricity (current as of 2024). Electrolysis keeps about
$65$ units, compression to $70\,\si{\mega\pascal}$ costs another few,
storage and handling a little more, and the fuel cell recovers about $55\%$
of what reaches it, leaving roughly $31$ units as electricity at the far
end. The losses compound multiplicatively, so the round trip returns about
a third of the input. The two large drops, at the electrolyser and at the
fuel cell, are where the engineering effort to raise the round trip
concentrates. The cascade is why hydrogen is reserved for long-duration
and hard-to-electrify uses where the round-trip penalty is paid rarely.}
\label{fig:vol04ch14:hydrogen-chain}
\end{figure}
